The thesis aims to develop and evaluate a hybrid multi-agent system for algorithmic trading, in which each agent relies on a distinct artificial intelligence technique. The system will utilize core financial inputs such as OHLCV price data, technical trading indicators (e.g., SMA, EMA, RSI, MACD), and sentiment signals derived from financial news, social media, and market sentiment indexes.

\vspace{3mm}
 
The multi-agent architecture will include forecasting agents based on selected time-series deep learning models (such as GRU, TCN, or Transformers), agents built using evolutionary algorithms for generating or optimizing trading rules, reinforcement learning agents interacting with a market simulation, as well as sentiment-driven NLP agents and simple rule-based or random baselines. A supervisory agent will aggregate signals from all underlying agents and determine the final trading decision, taking into account agent performance and user-defined preferences such as risk tolerance.

\vspace{3mm}

The research will compare the performance, robustness, and risk-adjusted results of individual agents and the overall system. The goal is to assess whether combining heterogeneous AI methods can yield more stable and reliable trading signals than any single approach.


